%% ma: L10-B14-0090
%% lop: 10
%% chuong: 5
%% bai: 14
%% dang: 10.14.9
%% loai: DS
%% muc: [NB, TH, TH, VD]
%% dap_an: "ĐSĐĐ"
%% nguon: "chuyen-de-cac-so-dac-trung-cua-mau-so-lieu-khong-ghep-nhom-toan-10-chuong-trinh-moi.pdf (GV. Phan Nhật Linh) (P05), Bài 03 Dạng 2 Phần II Câu 7"
%% kien_thuc: "Lớp 10 Bài 13 (số trung bình, trung vị, tứ phân vị, mốt), Bài 14 (khoảng biến thiên, khoảng tứ phân vị, phương sai, độ lệch chuẩn, giá trị bất thường)"
%% trang_thai: chinh-thuc
%% ghi_chu: "Cùng số liệu Bài tập tự luận 5 (đổi tên) nhưng khác loại câu. Tự duyệt theo yêu cầu giáo viên 22/10/2026."
\begin{ex}
Dưới đây là điểm kiểm tra giữa kì I của hai bạn An và Bình:
\begin{center}
\begin{tabular}{|c|c|c|c|c|c|c|c|c|c|c|c|}\hline
 & Toán & Lí & Hóa & Sinh & Anh & Văn & Sử & Địa & GDCD & Tin & CN\\\hline
An & $8$ & $7{,}5$ & $7{,}8$ & $8{,}3$ & $9$ & $7$ & $8$ & $8{,}2$ & $9$ & $8$ & $8{,}3$\\\hline
Bình & $8{,}5$ & $9{,}5$ & $9{,}5$ & $8{,}5$ & $9$ & $5$ & $5{,}5$ & $6$ & $10$ & $9$ & $8{,}5$\\\hline
\end{tabular}
\end{center}
Xét tính đúng sai của các khẳng định sau:
\choiceTF{\True Khoảng biến thiên của mẫu số liệu điểm kiểm tra của An là $2$}{Điểm kiểm tra trung bình của An là $8{,}3$}{\True Điểm kiểm tra trung bình của Bình (làm tròn đến hàng phần mười) là $8{,}1$}{\True Bạn An học đều hơn bạn Bình}
\loigiai{a) Đúng: $R=9-7=2$.

b) Sai: $\overline x_{An}=\dfrac{89{,}1}{11}=8{,}1$.

c) Đúng: $\overline x_{B}=\dfrac{89}{11}\approx8{,}1$.

d) Đúng: $s^2_{An}\approx0{,}31<s^2_{B}\approx2{,}76$.}
\end{ex}
